\documentclass[11pt,a4paper]{article}

\usepackage[T1]{fontenc}
\usepackage[utf8]{inputenc}
\usepackage{lmodern}
\usepackage{microtype}
\usepackage[margin=2.4cm]{geometry}
\usepackage{booktabs}
\usepackage{xcolor}
\usepackage{listings}
\usepackage{amsmath}
\usepackage{amssymb}
\usepackage[hidelinks]{hyperref}

\definecolor{codebg}{gray}{0.96}
\lstset{
  basicstyle=\ttfamily\footnotesize,
  backgroundcolor=\color{codebg},
  frame=none,
  columns=fullflexible,
  breaklines=true,
  showstringspaces=false,
  language=C,
  keywordstyle=\color{black}\bfseries,
  commentstyle=\itshape\color{gray!140},
  xleftmargin=1em,
  framexleftmargin=1em,
}

\newcommand{\code}[1]{\texttt{\small #1}}

\title{\textbf{Speeding up plain local search KNN graph construction}\\[2pt]
\large A profile-driven optimisation of the SISAP 2026 entry}
\author{Report on changes to \code{sisap2026-plain-local-search}}
\date{24 July 2026}

\begin{document}
\maketitle

\begin{abstract}
\noindent
The existing implementation builds a 15-nearest-neighbour graph over
$6\,350\,000$ Wikipedia \textsc{bge-m3} embeddings in \textbf{191.8\,s} on eight
threads, reaching a final recall of \textbf{0.8007}. After the changes described
here it does the same work in \textbf{118.1\,s} at recall \textbf{0.8009}: a
speedup of \textbf{1.62$\times$} with no loss of quality. The candidate recall
after every individual stage is statistically identical to the original, so this
is the same algorithm executed faster rather than a different quality/time
trade-off. Every change was driven by a measurement, and several plausible
optimisations were implemented, measured, found wanting, and reverted; those are
documented too, together with a measurement of the accuracy ceiling that the
1-bit quantisation imposes on the whole approach.
\end{abstract}

\section{Result}

\begin{table}[h]
\centering
\begin{tabular}{lrrr}
\toprule
Phase & Before & After & Change \\
\midrule
1-bit quantisation            & 0.2\,s   & 0.2\,s   & --- \\
Spatial relabelling + rerank codes & ---      & 2.0\,s   & new \\
Tree indexing                 & 64.2\,s  & 45.2\,s  & $-19.0$\,s \\
\quad \emph{tree partitioning}   & \emph{---} & \emph{12.9\,s} & \\
\quad \emph{leaf all-pairs}      & \emph{---} & \emph{29.9\,s} & \\
Refinement (2 rounds)         & 99.8\,s  & 48.1\,s  & $-51.7$\,s \\
\quad \emph{round 1}             & \emph{69.3\,s} & \emph{32.8\,s} & \\
\quad \emph{round 2}             & \emph{30.5\,s} & \emph{15.7\,s} & \\
Reranking                     & 27.0\,s  & 22.6\,s  & $-4.4$\,s \\
\midrule
\textbf{Total (build + query)} & \textbf{191.8\,s} & \textbf{118.1\,s} & $\mathbf{1.62\times}$ \\
\textbf{Final recall}          & \textbf{0.8007}   & \textbf{0.8009}   & \\
\bottomrule
\end{tabular}
\caption{Wall-clock per phase, \code{OMP\_NUM\_THREADS=8}. Times exclude HDF5
loading (which is outside the measured region in both versions) and exclude the
ground-truth recall evaluations, which are now subtracted explicitly.}
\end{table}

\noindent
Candidate recall (the fraction of true neighbours present anywhere in the
64-entry candidate list) matches the original at every stage:

\begin{center}
\begin{tabular}{lcc}
\toprule
After & Before & After changes \\
\midrule
indexing        & 0.2097 & 0.2082 \\
refinement r.1  & 0.7623 & 0.7621 \\
refinement r.2  & 0.7860 & 0.7861 \\
\midrule
final (reranked, top-15) & 0.8007 & 0.8009 \\
\bottomrule
\end{tabular}
\end{center}

\noindent
The small differences are run-to-run noise: repeated runs of the unmodified
binary varied over 0.76213--0.76230 for round~1 alone, because
insertion order across threads is nondeterministic.

\section{Setup and method}

All measurements were taken on an AMD EPYC 9754 (Zen 4c, 128 cores, 755\,GB
RAM, 16\,MB L3 per core-complex), compiled with the project's existing flags
(\code{-O3 -march=x86-64-v3 -fopenmp}) and run with exactly eight OpenMP
threads. The dataset is \code{benchmark-dev-wikipedia-bge-m3.h5}:
$m = 6\,350\,000$ vectors of $n = 1024$ half-precision floats, with a
$6\,350\,000 \times 32$ exhaustive ground truth under \code{/allknn/knns}.

Binary codes are $1024$ bits, so $bn = 16$ 64-bit words, exactly $128$ bytes ---
two cache lines --- per point. The code matrix is $813$\,MB, the neighbour
tables $V$ and $D$ are $1.6$\,GB each, and the fp16 input is $13$\,GB. These
sizes matter: almost every hot loop in this program is a random walk over one of
them, and the analysis below is mostly about cache lines rather than
instructions.

Two tools were built alongside the changes and are what most of the reasoning
rests on.

\paragraph{A kernel microbenchmark.} A single-threaded harness that replays one
leaf of the indexing phase ($L = 190$ points, all $17\,955$ pairs) in isolation,
against both an empty neighbour graph and a saturated one. Isolating this loop
made it possible to iterate on the insert path in seconds rather than in
four-minute full runs, and it immediately falsified the first hypothesis about
where indexing time was going.

\paragraph{An independent verifier.} A separate program that opens the produced
output file and the ground-truth file, and recomputes recall in the
\emph{original} index space, checking additionally for self-loops, duplicate
entries within a row, empty slots and out-of-range identifiers. Since one of the
changes relabels every point, the in-process recall number alone would not have
caught a mistake in the relabelling round-trip.

\section{Where the time actually went}

The original profile is dominated by three things: indexing (64.2\,s), the two
refinement rounds (99.8\,s), and reranking (27.0\,s). My first instinct was that
indexing was bound by lock contention and by the \code{memmove} inside the
sorted insert. A restructuring that removed \emph{all} locks from the leaf phase
returned indexing at 62.4\,s --- a 3\% improvement. The model was wrong.

The microbenchmark settled it:

\begin{center}
\begin{tabular}{lrr}
\toprule
& ns / pair & cycles @ 2.4\,GHz \\
\midrule
\code{hamming\_avx2\_lut} alone            & 4.76 & 11.4 \\
Full leaf loop, saturated graph            & 6.45 & 15.5 \\
Full leaf loop, empty graph                & 53.16 & 127.6 \\
\bottomrule
\end{tabular}
\end{center}

\noindent
The distance kernel is not the problem, and neither are locks. When the
neighbourhoods are still filling, nearly every candidate is accepted, and the
cost of \emph{landing} an accepted candidate is roughly eight times the cost of
computing the distance that justified it. Real indexing time sits between the
two regimes (13.5\,ns/pair measured across the whole phase), because early trees
run cold and later trees run warm.

The same reframing applies to refinement, where each candidate cost 238 cycles
against an 11.4-cycle distance --- the loop was waiting on memory, not
computing.

\section{Changes}

\subsection{Indexing: 64.2\,s $\rightarrow$ 45.2\,s}

\paragraph{Lock-free leaves.} Tree construction and the leaf all-pairs work are
now separate phases. The original ran whole trees concurrently, so every insert
went through the 8192-entry lock table. But all leaves of a \emph{single} tree
are disjoint point sets, so if one tree is processed at a time with its leaves
in parallel, there are no write conflicts at all and the leaf phase needs no
synchronisation whatsoever. Partitioning is parallelised with OpenMP tasks,
which also removes the previous serial bottleneck at the top of each tree.

\paragraph{An insert built for the indexing regime.} Since most offers succeed
here, the sorted insert's \code{memmove} is the dominant cost. During indexing
the neighbourhood is now kept \emph{unsorted} with the current maximum pinned to
slot~0; landing a candidate is a slot-0 overwrite plus a rescan for the new
maximum, and the ordering is restored once, afterwards, by a single pass
(\code{KNN\_sort\_lists}). Two further details made this actually fast:

\begin{itemize}
\item \textbf{No store-forward stall.} Writing \code{D[0]} and then immediately
reloading it as part of a 32-byte vector incurs a store-forwarding failure on
every successful offer. Slot~0 is instead \emph{blended} with the incoming
distance in registers, so nothing is stored before the reduction.
\item \textbf{No argmax scan.} Locating the new maximum needs its position, not
just its value. Rather than a second, badly-predicted search loop, each distance
carries its own slot index in the low \code{KNN\_TAG\_BITS} bits, so a single
\code{vpmaxsd} reduction yields both the value and where it lives.
\end{itemize}

\noindent
The first two took the empty-graph leaf cost from 53.2 to 27.6\,ns/pair; storing
the tags in \code{D} rather than recomputing them in the reduction loop brought it
to 25.6\,ns/pair, an overall $2.1\times$, with the saturated case unchanged. Note
how little the last step gave: by then the offer is at its load-throughput floor,
which is why the insert path was not pursued further.

\paragraph{Prefetching the partition scan.} The pivot-distance pass reads
\code{Matrix[O[i].v]} for scattered \code{v}, i.e.\ a random 128-byte read out
of 813\,MB, and all the identifiers are known ahead of time. Running the loads
twelve iterations ahead cut partitioning from 19.8\,s to 12.6\,s.

\paragraph{Leaf-local scratch buffer.} A leaf's codes are gathered into a
contiguous $\approx$25\,KB buffer before the all-pairs loop, so the inner loop
streams over an L1-resident block instead of chasing 190 rows spread over
813\,MB. The gather itself is prefetched.

\paragraph{Bounded histogram.} The median split histogram is cleared and scanned
only over the observed distance range. In 1024 dimensions the distances to a
random pivot concentrate very tightly, so this is a narrow window out of the
full $[0, 1024]$ span. (The original also indexed \code{C[d]} with $d$ up to
1024 into an array of size $bn \cdot 64 = 1024$, a one-element overrun; the new
histogram is sized 1088.)

\subsection{Refinement: 99.8\,s $\rightarrow$ 48.1\,s}

\paragraph{Candidate deduplication --- the largest single win.} Nearby points
share neighbours, so the same $w$ is reached through many different $v$ while
scanning $u$'s neighbours-of-neighbours. Each repeat costs a full distance
evaluation and three random cache lines. A small per-thread direct-mapped hash
table (2048 entries, 16\,KB, keyed on $(u,w)$) drops the repeats for a few
cycles each.

This is worth stating carefully, because skipping work in a search loop usually
does change the answer. Here it does not:

\begin{quote}
\emph{Claim.} Suppressing a repeated offer of $(u,w)$ leaves the final graph
unchanged.

\emph{Argument.} The pair $(u,w)$ has one fixed distance $d$. \code{Dmax[u]} is
non-increasing over time. If the first offer was rejected, $d \geq
\code{Dmax}[u]$ then, and \code{Dmax}[u] is no larger later, so any repeat is
rejected too. If it was accepted, $w$ is in $u$'s list and the insert's
duplicate check makes any repeat a no-op. The same argument applies to the
reverse insert into $w$. Hash collisions can only cause a \emph{false miss} ---
the key comparison is exact, so a colliding entry compares unequal and the
candidate is evaluated normally.
\end{quote}

\noindent
The measured candidate recall was indeed unchanged.

\paragraph{A compact \code{Dmax} array.} Rejecting a candidate is the common
case and requires the worst distance in $w$'s neighbourhood. The original read
\code{D[w*k + k-1]}, pulling a cache line out of a 1.6\,GB table. That value is
now mirrored in a compact $m$-element array, so the test touches four bytes out
of 25\,MB.

\paragraph{Prefetching the dependency chain.} The candidate identifiers come
from $v$'s neighbour list, which is itself a cold random read; nothing in an
iteration could be prefetched until it landed. The scan now runs the code rows
and \code{Dmax} entries eight candidates ahead, \emph{and} starts the next $v$'s
neighbour list at the top of each iteration. That last one alone was worth
5.2\,s in round~1.

\paragraph{Inlined shift.} The sorted insert opens a hole with \code{memmove}
of a runtime length, which compiles to a library call whose dispatch overhead
dominates a shift this short. It is now done inline in 8-integer AVX2 blocks
from the top down; each block's source is disjoint from the previous block's
destination, and the block at the insertion point is loaded before the loop
because the loop overwrites it (with identical values, so the final store stays
consistent).

\paragraph{A cheaper lock.} \code{omp\_lock\_t} costs an out-of-line call into
libgomp on both acquire and release. These locks are held for a few dozen cycles
and are essentially never contended (eight threads over 8192 locks), so they were
replaced with an inline test-and-set spinlock.

\subsection{Reranking: 27.0\,s $\rightarrow$ 22.6\,s}

Reranking computes $6.35\times10^6 \times 64 \approx 4.1\times10^8$ exact fp16
inner products, each reading a 2\,KB row at random from a 13\,GB table --- about
812\,GB of traffic. At 1140 cycles per dot product against roughly 200 cycles of
arithmetic, it is almost entirely memory-bound. It benefits from the spatial
relabelling below, from prefetching the next candidate's row, and from the same
compact-array trick as \code{Dmax}: the reverse-insert test now reads a 4-byte
entry from a 25\,MB \code{Rmin} array instead of a cache line from the 381\,MB
result table.

\subsection{Memory layout}

\paragraph{Spatial relabelling.} Points are relabelled along the leaf order of
one extra partition pass, so that points close in space are close in memory, and
everything downstream runs in the new labelling; only the output is mapped back.
The extra pass costs 0.6\,s and is worth 5.9\,s (measured directly with
\code{KNN\_NOPERM=1}: 123.8\,s versus 118.0\,s). It is a real but modest win ---
a single randomised vantage-point tree gives good locality only near the leaves,
and a stronger ordering (for example reordering by the neighbour graph itself
after indexing) is the obvious place to look for more.

\paragraph{Huge pages.} All large allocations are 2\,MB-aligned and marked
\code{MADV\_HUGEPAGE}. With 4\,KB pages, random access over these arrays misses
the TLB on essentially every access. This has \emph{no effect on the test
machine}: \code{transparent\_hugepage/enabled} reads \code{[madvise]}, but no
huge pages are actually granted --- \code{thp\_fault\_alloc} does not move for a
1\,GB \code{mmap} plus \code{madvise}, and there is no \code{hugetlb} pool. The
calls are free and would help on a host where THP works, so they were kept.

\section{Measured and rejected}

Several changes looked good on paper, were implemented and measured, and did not
survive. They are listed because the negative results constrain what is worth
trying next.

\begin{table}[h]
\centering
\begin{tabular}{lrrrr}
\toprule
Variant (vs.\ its own reference run) & \multicolumn{2}{c}{variant} & \multicolumn{2}{c}{reference} \\
\cmidrule(lr){2-3}\cmidrule(lr){4-5}
 & total & recall & total & recall \\
\midrule
int8 rerank codes                & 157.9\,s & 0.7978 & 160.6\,s & 0.8007 \\
Rerank only top 48 of 64         & 114.3\,s & 0.7796 & 119.2\,s & 0.8006 \\
Candidates scored in index order & 118.0\,s & 0.8008 & 118.0\,s & 0.8008 \\
No spatial relabelling           & 123.8\,s & 0.8007 & 118.0\,s & 0.8008 \\
\bottomrule
\end{tabular}
\caption{Rejected variants, each against the run it was A/B-tested against ---
they were measured at different points during the work, so their absolute totals
are not directly comparable to each other. Recall differences are well outside
run-to-run noise.}
\end{table}

\paragraph{int8 rerank codes.} Halving the bytes per vector should have helped a
memory-bound loop, and it did --- but only by 2.7\,s, because the loop is limited
by page walks and access latency more than by bandwidth. It cost 0.0029 recall,
which at this operating point is not affordable. (The exact fp16 path is kept;
the int8 one is still reachable via \code{KNN\_RERANK\_F16=0}.)

\paragraph{Truncating the rerank.} Reranking only the best 48 candidates by
Hamming distance saves 5\,s but costs 0.021 recall. The tail of the candidate
list matters more than its own ranking suggests, because reranking also inserts
$u$ into each candidate's result list --- those reverse edges are what lift the
final recall (0.8008) above the candidate recall (0.7861) in the first place.

\paragraph{Rebalancing indexing against refinement.} Indexing consumes $5\times$
more distance computations than both refinement rounds combined, yet only
reaches candidate recall 0.209, which invites the conclusion that it is
over-provisioned. It is not. Refinement \emph{saturates}, and the plateau it
saturates at is set by the quality of the initial graph:

\begin{table}[h]
\centering
\begin{tabular}{lrrrrrr}
\toprule
Trees / rounds & Indexing & \multicolumn{3}{c}{candidate recall after round} & Total & Final \\
\cmidrule(lr){3-5}
 & & 1 & 2 & 3 (4) & & recall \\
\midrule
 8 / 4 & 17.8\,s & 0.6990 & 0.7432 & 0.7509 (0.7527) & 131.4\,s & 0.7730 \\
16 / 3 & 28.2\,s & 0.7318 & 0.7648 & 0.7704          & 121.1\,s & 0.7879 \\
24 / 2 & 37.8\,s & 0.7490 & 0.7767 & ---             & 112.8\,s & 0.7920 \\
\textbf{32 / 2} & 46.9\,s & 0.7620 & 0.7859 & ---    & \textbf{119.2\,s} & \textbf{0.8006} \\
\bottomrule
\end{tabular}
\caption{Fewer trees plus more refinement rounds is worse on \emph{both} axes.
A fourth round buys 0.0018 recall for 15\,s.}
\end{table}

\noindent
Every cheaper-indexing configuration lands below 0.80, and the extra rounds
needed to compensate cost more than the trees saved. The existing
$Q_b = 264$, $Q_n = 32$, $L_n = 2$ is a well-chosen operating point and was left
alone.

\section{How much headroom is there?}

A useful question when a target is only just being met: how much of the shortfall
is search quality, and how much is the representation? The candidate set is
chosen entirely in the 1-bit Hamming space, so anything the binary codes cannot
separate is lost regardless of how good the search is.

This is directly measurable. Brute-forcing the \emph{exact} top-64 by Hamming
distance for 300 sampled points and scoring it against the ground truth
(available as \code{KNN\_CEILING=300}) gives:

\begin{center}
\fbox{\parbox{0.86\textwidth}{\centering
Exact top-64 by Hamming distance contains \textbf{87.3\%} of the true
nearest neighbours.\\[2pt]
The pipeline reaches candidate recall \textbf{0.786} --- \textbf{90\%} of what
is achievable.}}
\end{center}

\noindent
So 1-bit quantisation at $k = 64$ caps the whole approach at roughly 0.87, the
search is already extracting 90\% of that, and the working margin above the 0.80
target is thin. This explains the failures in the previous section: there simply
is not enough slack to give any away. It also says where the remaining
opportunities are --- a richer code or a larger $k$ would raise the ceiling, at
a cost that would have to be earned back elsewhere.

\section{Verification}

The in-process recall figure is computed in the relabelled index space, so it
cannot by itself confirm that the relabelling is inverted correctly on output.
The output file was therefore checked independently, reading both the produced
\code{knns} dataset and \code{/allknn/knns} from the benchmark file and scoring
in the original labelling:

\begin{lstlisting}[language=]
answer 6350000 x 15, ground truth 6350000 x 32
recall            = 0.800851
self loops        = 0
empty slots       = 0
out of range      = 0
duplicates in row = 0
\end{lstlisting}

\noindent
In addition, a standalone unit test drives \code{KNN\_offer\_max0} and
\code{KNN\_insert\_unsafe} through 200 randomised trials, asserting after each
that slot~0 holds the maximum, that \code{Dmax} agrees with the list, that no
duplicates exist, that stored distances match the real pair distances, that the
list sorts correctly, and that offering the complete candidate set converges to
the exact top-$k$. This is what gave confidence in the tagged-distance
representation and the hand-written AVX2 shift, both of which are easy to get
subtly wrong.

\section{Notes for reuse}

\begin{itemize}
\item The two \code{if (recall > 0.8) goto endwork} early exits were removed.
They consult the ground truth, which will not be present on a real run, so they
made the measured configuration differ from the deployed one.

\item Parameters are overridable from the environment for experimentation;
defaults are the tuned configuration. \code{KNN\_QB}, \code{KNN\_QN},
\code{KNN\_K}, \code{KNN\_LN} are the original four; \code{KNN\_RR} limits how
many candidates are reranked, \code{KNN\_RERANK\_F16} switches between exact
fp16 and int8 codes, \code{KNN\_NOPERM} disables the relabelling,
\code{KNN\_CEILING=$n$} runs the quantisation-ceiling diagnostic and exits.

\item \code{KNN\_CACHE=<path>} is a development convenience only: a raw cache of
the parsed input so that repeated runs skip re-parsing 14.8\,GB of HDF5 over
NFS. Loading is outside the measured region either way, so it does not affect
any number in this report.

\item The build stays on \code{-march=x86-64-v3}. Worth noting for the future:
the test machine's Zen 4c cores support AVX-512 \textsc{vpopcntdq}, which would
cut the Hamming kernel by roughly $3\times$; behind a runtime
\code{\_\_builtin\_cpu\_supports} check this would remain safe on a v3-only
host. It was left out deliberately.

\item \code{compute\_dot\_product} in \code{matrix.h} is now unused, superseded
by an aligned \code{dot\_f16\_avx2}. It was left in place rather than modify a
file that otherwise needed no changes.
\end{itemize}

\section*{Appendix: incremental measurements}

Each row is a full run; the change column names what was added relative to the
row above.

\begin{center}
\begin{tabular}{llrr}
\toprule
\# & Change & Total & Final recall \\
\midrule
0 & original                                                & 191.8\,s & 0.8007 \\
1 & lock-free leaves, relabelling, int8 rerank              & 157.9\,s & 0.7978 \\
2 & exact fp16 rerank restored                              & 160.6\,s & 0.8007 \\
3 & blend + packed argmax, partition/gather prefetch        & 143.6\,s & 0.8009 \\
4 & candidate dedup, spinlock, inlined shift                & 125.1\,s & 0.8008 \\
5 & prefetch next neighbour list                            & 119.2\,s & 0.8006 \\
6 & pre-tagged distances                                    & 118.0\,s & 0.8008 \\
7 & confirmation run                                        & 118.1\,s & 0.8009 \\
\bottomrule
\end{tabular}
\end{center}

\vspace{1em}
\noindent
Files changed: \code{include/algorithms.h}, \code{include/knn.h},
\code{src/knn.c}, \code{src/main.c}. No commits were made; the working tree is
left dirty on top of \code{5011cae}.

\end{document}
